\documentclass[11pt,a4paper]{article}
\usepackage[margin=1in]{geometry}
\usepackage{amsmath,amssymb,amsfonts,amsthm}
\usepackage{booktabs}
\usepackage{hyperref}
\usepackage{xcolor}
\usepackage{tcolorbox}

\newtheorem{theorem}{Theorem}
\newtheorem{proposition}{Proposition}
\newtheorem{lemma}{Lemma}
\newtheorem{corollary}{Corollary}
\newtheorem{finding}{Finding}
\theoremstyle{definition}
\newtheorem{definition}{Definition}
\theoremstyle{remark}
\newtheorem{remark}{Remark}

\title{Derivation of $\kappa$ from the DFD Action Principle:\\
Does the Optical Metric Force $\kappa = 0$?}

\author{Derived for G.\ Alcock --- DFD Research Campaign\\
March 23, 2026}
\date{}

\begin{document}
\maketitle

%=============================================================================
\section{Executive Summary}
%=============================================================================

\begin{tcolorbox}[colback=red!5!white, colframe=red!60!black, title=Principal Finding]
\textbf{The optical metric forces $\kappa = 0$ at tree level.}
The EM Lagrangian on the DFD optical metric $\tilde{g}_{\mu\nu}$ treats
$E$ and $B$ symmetrically: $\varepsilon_{\rm eff} = \varepsilon_0 e^{+\psi}$,
$\mu_{\rm eff} = \mu_0 e^{+\psi}$ (the Tamm--Plebanski relations).  There
is no room for a differential $e^{\pm\kappa\psi}$ split.

\medskip
\textbf{Any nonzero $\kappa$ requires physics beyond the minimal DFD
postulates.}  The auxiliary covariant metric
$\hat{g}_{\mu\nu} = \mathrm{diag}(-c^2 e^{-2\psi},\; e^{2\psi}\delta_{ij})$
used in the gauge emergence derivations \textit{does} produce an E--B
asymmetry ($e^{+2\psi}$ vs $e^{-2\psi}$), and this is the origin of the
$\kappa \neq 0$ phenomenology.  But this auxiliary metric is a computational
device, not the physical metric.

\medskip
\textbf{Resolution:}  The $\kappa$ parameter measures the departure of the
\textit{gauge-sector effective metric} from the \textit{physical optical metric}.
This departure is generated by the gauge-$\psi$ coupling in the gauge emergence
framework.  From the auxiliary metric, $\kappa = 2$ at tree level.  Whether
this survives as the physical value or is screened to a smaller effective $\kappa$
is an open question.
\end{tcolorbox}

%=============================================================================
\section{The Three Metrics and Their EM Lagrangians}
%=============================================================================

DFD employs three related metrics.  Each gives a different EM Lagrangian
and a different value of $\kappa$.

\subsection{Metric 1: The Gordon Optical Interval ($\kappa = 0$)}

The fundamental DFD description uses the Gordon-type optical interval
(Section~2.1 of the main text):
\begin{equation}\label{eq:gordon}
d\tilde{s}^2_{\rm Gordon} = -\frac{c^2}{n^2}\,dt^2 + d\mathbf{x}^2,
\qquad n = e^\psi.
\end{equation}
The spatial sections are \textbf{flat Euclidean}.  Through the Tamm--Plebanski
construction (Section~10 of the main text), this defines effective vacuum
constitutive relations:
\begin{equation}\label{eq:tamm}
\varepsilon_{\rm eff} = \varepsilon_0\, n = \varepsilon_0\, e^{+\psi},
\qquad
\mu_{\rm eff} = \mu_0\, n = \mu_0\, e^{+\psi}.
\end{equation}
Both $\varepsilon$ and $\mu$ scale identically with $\psi$.  The impedance is
\begin{equation}
Z(\psi) = \sqrt{\frac{\mu_{\rm eff}}{\varepsilon_{\rm eff}}}
= \sqrt{\frac{\mu_0}{\varepsilon_0}} = Z_0,
\end{equation}
independent of $\psi$.  The EM energy density is
\begin{equation}\label{eq:uEM-gordon}
u_{\rm EM}^{\rm (Gordon)} = \frac{1}{2}\varepsilon_0 e^{+\psi} E^2
+ \frac{1}{2\mu_0} e^{+\psi} B^2
= \frac{e^{+\psi}}{2}\!\left(\varepsilon_0 E^2 + \frac{B^2}{\mu_0}\right).
\end{equation}
This is the $\kappa = 0$ case: $E$ and $B$ sectors couple to $\psi$
identically (both get $e^{+\psi}$).

\begin{finding}[Optical Metric Forces $\kappa = 0$]
The constitutive relations derived from the Gordon optical interval
(\ref{eq:gordon}) via the Tamm--Plebanski map are impedance-matched:
$\varepsilon_{\rm eff}/\varepsilon_0 = \mu_{\rm eff}/\mu_0 = n$.
This forces $\kappa = 0$ identically.
\end{finding}

\subsection{Metric 2: The Physical Metric ($\kappa = 0$)}

The physical metric through which matter couples (Section~2.2):
\begin{equation}\label{eq:physical}
\tilde{g}_{\mu\nu}^{\rm (phys)} = \mathrm{diag}(-c^2 e^{-\psi},\;
e^{+\psi},\; e^{+\psi},\; e^{+\psi}).
\end{equation}
Computing $\sqrt{-\tilde{g}} = c\, e^{+\psi}$ and the EM Lagrangian
$\mathcal{L}_{\rm EM} = -\frac{1}{4}\sqrt{-\tilde{g}}\,
\tilde{g}^{\mu\alpha}\tilde{g}^{\nu\beta} F_{\mu\nu}F_{\alpha\beta}$:
\begin{align}
\tilde{g}^{00} &= -e^{+\psi}/c^2, \qquad
\tilde{g}^{ij} = e^{-\psi}\,\delta^{ij}, \\[4pt]
\tilde{g}^{0\alpha}\tilde{g}^{0\beta}F_{0i}F_{0j}\delta^{ij}
&\propto e^{+\psi}\cdot e^{-\psi}\cdot E^2 = E^2, \\[4pt]
\tilde{g}^{i\alpha}\tilde{g}^{j\beta}F_{ij}F_{\alpha\beta}
&\propto e^{-\psi}\cdot e^{-\psi}\cdot B^2 = e^{-2\psi}\,B^2.
\end{align}
With the $\sqrt{-\tilde{g}} = c\,e^{+\psi}$ prefactor:
\begin{equation}\label{eq:LEM-physical}
\mathcal{L}_{\rm EM}^{\rm (phys)} \propto
e^{+\psi}\,E^2 - e^{-\psi}\,B^2.
\end{equation}
Wait---this gives \textit{different} exponents for $E$ and $B$!
Identifying with the constitutive form
$u_{\rm EM} = \frac{1}{2}\varepsilon_0 n\,e^{+\kappa\psi}\,E^2
+ \frac{1}{2\mu_0}n\,e^{-\kappa\psi}\,B^2$
where $n = e^\psi$:
\begin{itemize}
\item $E$-sector exponent: $(1+\kappa)\psi$.  From (\ref{eq:LEM-physical}): $+\psi$.
  So $1 + \kappa = 1$, giving $\kappa = 0$.
\item $B$-sector exponent: $(1-\kappa)\psi$.  From (\ref{eq:LEM-physical}): $-\psi$.
  So $1 - \kappa = -1$, giving $\kappa = 2$.
\end{itemize}

\begin{remark}[Apparent Inconsistency]
The physical metric gives conflicting values: $\kappa = 0$ from the
$E$-sector and $\kappa = 2$ from the $B$-sector.  This is resolved by
noting that the constitutive parametrization
$\varepsilon = \varepsilon_0\,n\,e^{+\kappa\psi}$,
$\mu = \mu_0\,n\,e^{-\kappa\psi}$ assumes the \textit{product}
$\varepsilon\mu = \varepsilon_0\mu_0\,n^2$ is preserved.  The physical
metric does preserve $v_{\rm ph} = c/n$, but the E-B split is
$e^{+\psi}$ vs $e^{-\psi}$, which corresponds to the parametrization
\begin{equation}
\varepsilon = \varepsilon_0\,e^{+\psi}, \qquad
\mu = \mu_0\,e^{-\psi}, \qquad
\varepsilon\mu = \varepsilon_0\mu_0.
\end{equation}
This is a \textbf{different parametrization} from the $\kappa$ convention.
In the $\kappa$ convention (where $\varepsilon\mu = \varepsilon_0\mu_0 n^2$),
we must factor out $n = e^\psi$ first, giving the residual split
$e^{0} = 1$ for both sectors: $\kappa = 0$.

The resolution is that the physical metric gives $\kappa = 0$ in the
standard convention because the asymmetry in the metric is entirely
accounted for by the $n = e^\psi$ prefactor.
\end{remark}

\subsection{Metric 3: The Auxiliary Covariant Metric ($\kappa_{\rm eff} = 2$)}

The auxiliary metric used in gauge emergence calculations (Appendix~G):
\begin{equation}\label{eq:auxiliary}
\hat{g}_{\mu\nu} = \mathrm{diag}(-c^2 e^{-2\psi},\;
e^{+2\psi},\; e^{+2\psi},\; e^{+2\psi}).
\end{equation}
This ``doubles the exponents'' relative to the physical metric.  The
EM Lagrangian on this metric (as derived in Appendix~G) gives:
\begin{equation}\label{eq:LEM-aux}
\mathcal{L}_{\rm YM}^{(r)} = \frac{e^{+2\psi}}{2g_r^2 c}\,E_r^2
- \frac{c\,e^{-2\psi}}{2g_r^2}\,B_r^2.
\end{equation}
Now identifying with the $\kappa$ parametrization.  Write
$n_{\rm aux} = e^{2\psi}$ (the ``refractive index'' of the auxiliary
metric).  Then:
\begin{itemize}
\item $E$-sector: $e^{+2\psi} = n_{\rm aux}\,e^{+\kappa_{\rm aux}\psi}$
  $= e^{2\psi}\,e^{+\kappa_{\rm aux}\psi}$.
  So $\kappa_{\rm aux} = 0$ from the $E$-sector.
\item $B$-sector: $e^{-2\psi} = n_{\rm aux}\,e^{-\kappa_{\rm aux}\psi}$
  $= e^{2\psi}\,e^{-\kappa_{\rm aux}\psi}$, needing $e^{-2\psi} = e^{(2-\kappa_{\rm aux})\psi}$.
  So $\kappa_{\rm aux} = 4$.
\end{itemize}

This is again in the wrong convention.  In the standard $\kappa$ convention
(where $v_{\rm ph} = c/n$ with $n = e^\psi$, and
$\varepsilon\mu = \varepsilon_0\mu_0 n^2$):
\begin{equation}
\varepsilon_{\rm aux} \propto e^{+2\psi}, \qquad
\mu_{\rm aux} \propto e^{-2\psi}, \qquad
\varepsilon_{\rm aux}\mu_{\rm aux} \propto 1 = \varepsilon_0\mu_0\,n^0.
\end{equation}
Factoring out $n = e^\psi$:
\begin{equation}
\varepsilon_{\rm aux} = \varepsilon_0\,n\,e^{+\psi},
\qquad
\mu_{\rm aux} = \mu_0\,n^{-1}\,e^{-3\psi}.
\end{equation}
This does not match the standard $\kappa$ parametrization cleanly.

\begin{finding}[The Auxiliary Metric Generates a Real E--B Asymmetry]
The auxiliary metric $\hat{g}_{\mu\nu}$ produces the Lagrangian
$\mathcal{L} \propto e^{+2\psi}E^2 - e^{-2\psi}B^2$, which has a
manifest E--B asymmetry under $\psi$-variation:
\begin{equation}
\frac{\partial\mathcal{L}}{\partial\psi}
\propto +2 e^{+2\psi} E^2 + 2 e^{-2\psi} B^2.
\end{equation}
Both terms source $\psi$ with the \textbf{same sign}, but with different
$\psi$-dependences.  This is the mathematical origin of the $k_a$
derivation in Appendix~G: in the magnetically dominated regime, the
$B^2$ term dominates, giving the $1/(4\alpha)$ factor.
\end{finding}

%=============================================================================
\section{Why the Optical Metric Forces $\kappa = 0$: The Proof}
%=============================================================================

\begin{theorem}[Optical Metric $\Rightarrow$ $\kappa = 0$]\label{thm:kappa-zero}
Let $\psi$ be the DFD scalar field and let the EM sector propagate on
the Gordon optical interval $d\tilde{s}^2 = -c^2 n^{-2}dt^2 + d\mathbf{x}^2$
with $n = e^\psi$.  Then the effective constitutive relations are
impedance-matched ($Z = Z_0$), and $\kappa = 0$ identically.
\end{theorem}

\begin{proof}
The Gordon optical metric has flat spatial sections ($g_{ij} = \delta_{ij}$)
and a position-dependent lapse $N = c/n$.

The Tamm--Plebanski construction for EM in a medium with metric
$g_{\mu\nu} = \mathrm{diag}(-N^2, \delta_{ij})$ gives:
\begin{equation}
\varepsilon_{\rm eff}^{ij} = \mu_{\rm eff}^{ij}
= \frac{\sqrt{\gamma}}{N}\,\gamma^{ij}
= \frac{n}{c}\,\delta^{ij}.
\end{equation}
Since $\gamma_{ij} = \delta_{ij}$ (flat spatial sections), we have
$\sqrt{\gamma} = 1$ and $\gamma^{ij} = \delta^{ij}$.  Thus
\begin{equation}
\varepsilon_{\rm eff} = \varepsilon_0\,n, \qquad
\mu_{\rm eff} = \mu_0\,n.
\end{equation}
Both scale identically with $n = e^\psi$.  The impedance
$Z = \sqrt{\mu_{\rm eff}/\varepsilon_{\rm eff}} = \sqrt{\mu_0/\varepsilon_0} = Z_0$
is independent of $\psi$.

In the $\kappa$ parametrization
$\varepsilon = \varepsilon_0\,n\,e^{+\kappa\psi}$,
$\mu = \mu_0\,n\,e^{-\kappa\psi}$, the impedance becomes
$Z = Z_0\,e^{-\kappa\psi}$.  Since $Z = Z_0$ for all $\psi$, we
conclude $\kappa = 0$.
\end{proof}

\begin{corollary}[Nonzero $\kappa$ Requires Beyond-Optical-Metric Physics]
Any $\kappa \neq 0$ cannot arise from the minimal DFD postulates
(optical metric + scalar field action).  It must originate from
additional structure in the theory.
\end{corollary}

%=============================================================================
\section{What Physics Could Generate $\kappa \neq 0$?}
%=============================================================================

We identify four potential sources of nonzero $\kappa$.

\subsection{Source 1: The Gauge Emergence Auxiliary Metric}

The gauge emergence framework (Appendix~G) uses the auxiliary metric
$\hat{g}_{\mu\nu} = \mathrm{diag}(-c^2 e^{-2\psi}, e^{2\psi}\delta_{ij})$,
which is a ``doubled-exponent'' version of the physical metric.  On this
metric, the Yang--Mills Lagrangian is
\begin{equation}
\mathcal{L}_{\rm YM} = \frac{e^{+2\psi}}{2g^2 c}\,E^2
- \frac{c\,e^{-2\psi}}{2g^2}\,B^2.
\end{equation}
The $\psi$-variation of this Lagrangian sources the $\psi$ field
asymmetrically in $E$ and $B$.  If the gauge fields ``see'' the
auxiliary metric rather than the Gordon metric, then $\kappa \neq 0$
is generated at tree level.

\paragraph{Physical mechanism.}
In the gauge emergence picture, gauge fields arise as Berry connections
on the internal space $\mathbb{C}P^2 \times S^3$.  The frame stiffness
$\kappa_r = n_r\,\kappa_0$ generates a metric for gauge fluctuations
that differs from the metric for matter or light propagation.  This
\textit{bi-metric} structure---physical metric for matter, auxiliary
metric for gauge dynamics---is the origin of $\kappa \neq 0$.

\paragraph{Effective $\kappa$ from the auxiliary metric.}
The auxiliary metric gives $\varepsilon \propto e^{+2\psi}$ and
$\mu^{-1} \propto e^{+2\psi}$, i.e., $\mu \propto e^{-2\psi}$.
In the standard $\kappa$ convention with $n = e^\psi$:
\begin{align}
\varepsilon &= \varepsilon_0\,n\,e^{+\kappa\psi}
= \varepsilon_0\,e^{(1+\kappa)\psi}
\stackrel{!}{=} \varepsilon_0\,e^{+2\psi}
\quad\Rightarrow\quad \kappa = 1, \\[4pt]
\mu &= \mu_0\,n\,e^{-\kappa\psi}
= \mu_0\,e^{(1-\kappa)\psi}
\stackrel{!}{=} \mu_0\,e^{-2\psi}
\quad\Rightarrow\quad \kappa = 3.
\end{align}
These are inconsistent because the auxiliary metric does not preserve
$\varepsilon\mu = \varepsilon_0\mu_0\,n^2$.  Instead,
$\varepsilon_{\rm aux}\,\mu_{\rm aux} = \varepsilon_0\mu_0$ (no $n^2$ factor).

This tells us the $\kappa$ parametrization (which assumes
$v_{\rm ph} = c/n$) does not directly apply to the auxiliary metric.
The auxiliary metric has its own ``effective refractive index''
$n_{\rm aux} = e^{2\psi}$ with $v_{\rm ph,aux} = c\,e^{-2\psi}$.
\textbf{The correct interpretation is that the gauge emergence
auxiliary metric is not a constitutive modification of the optical
metric; it is a genuinely different metric for gauge field dynamics.}

\subsection{Source 2: Conformal Anomaly / Trace Anomaly}

At one loop, the EM stress tensor acquires a trace anomaly:
\begin{equation}
\langle T^\mu{}_\mu \rangle = \frac{\alpha}{90\pi}\left(R_{\mu\nu\alpha\beta}^2
- R_{\mu\nu}^2 + \Box R\right).
\end{equation}
On the optical metric background, this generates a correction to the
$\varepsilon$--$\mu$ split because the Riemann tensor components for
the Gordon metric differ from those of a conformally flat metric.
This could generate an effective $\kappa$ at order $\alpha$.

\paragraph{Estimate.}
$\kappa_{\rm anomaly} \sim \alpha/(2\pi) \sim 10^{-3}$, much smaller
than unity.

\subsection{Source 3: Running Gauge Couplings in $\psi$}

If the gauge coupling ``runs'' with $\psi$ (through the dependence
of the renormalization scale on the local gravitational potential),
then $\alpha_{\rm eff}(\psi) = \alpha_0 + \delta\alpha\cdot\psi + \ldots$
This modifies the Coulomb interaction but not the magnetic interaction
at the same order, generating an effective $\kappa$:
\begin{equation}
\kappa_{\rm running} = \frac{\partial\ln\alpha}{\partial\psi}
= k_\alpha = \frac{\alpha^2}{2\pi} \approx 8.5 \times 10^{-6}.
\end{equation}
This is the $k_\alpha$ parameter derived in Appendix~G
(Theorem~G.5), and it is extremely small.

\subsection{Source 4: Disformal Coupling (Speculative)}

A disformal metric of the form
\begin{equation}
\bar{g}_{\mu\nu} = \tilde{g}_{\mu\nu}
+ D(\psi)\,\partial_\mu\psi\,\partial_\nu\psi
\end{equation}
would modify the EM Lagrangian in a way that breaks the E--B symmetry.
The disformal factor $D$ would introduce direction-dependent propagation
speeds, distinguishing electric (timelike) from magnetic (spacelike)
components of $F_{\mu\nu}$.

For a static field $\psi = \psi(\mathbf{x})$, the disformal correction
modifies only the spatial metric:
\begin{equation}
\bar{g}_{ij} = e^{+\psi}\,\delta_{ij}
+ D\,\partial_i\psi\,\partial_j\psi.
\end{equation}
This introduces an anisotropy rather than a uniform $\kappa$ split,
so it would not produce the simple parametrization in
Eq.~(\ref{eq:tamm}).  Disformal coupling is therefore
\textbf{not the origin of $\kappa$}.

%=============================================================================
\section{Definitive Analysis: What the Action Principle Says}
%=============================================================================

\subsection{The Full DFD + EM Action}

The complete action including EM is:
\begin{equation}
S = S_\psi + S_{\rm EM} + S_{\rm matter},
\end{equation}
where $S_\psi$ is given by Eq.~(5) of the main text and
\begin{equation}\label{eq:SEM}
S_{\rm EM} = -\frac{1}{4}\int d^4x\,\sqrt{-\tilde{g}}\;
\tilde{g}^{\mu\alpha}\tilde{g}^{\nu\beta}\,F_{\mu\nu}F_{\alpha\beta}.
\end{equation}
On the Gordon metric (\ref{eq:gordon}):
\begin{align}
\sqrt{-\tilde{g}} &= c/n = c\,e^{-\psi}, \\
\tilde{g}^{00} &= -n^2/c^2 = -e^{2\psi}/c^2, \qquad
\tilde{g}^{ij} = \delta^{ij}.
\end{align}
Computing the EM Lagrangian density:
\begin{align}
\mathcal{L}_{\rm EM}^{\rm (Gordon)}
&= -\frac{c\,e^{-\psi}}{4}\left[
  2\cdot\frac{e^{2\psi}}{c^2}\cdot\frac{E^2}{c^2}\cdot c^2
  - B^2\right] \notag\\
&= -\frac{c\,e^{-\psi}}{4}\left[
  \frac{2e^{2\psi}}{c^2}\,E^2 - B^2\right] \notag\\
&= \frac{e^{-\psi}}{2}\left[
  -\frac{e^{2\psi}}{c}\,E^2 + \frac{c}{2}\,B^2\right].
\end{align}

Let me redo this more carefully.  $F_{0i} = E_i/c$ (SI conventions
with $A_0$ having dimensions of voltage).  Actually, let us use
Gaussian-like conventions where $F_{0i} = E_i$ and
$F_{ij} = \epsilon_{ijk}B_k$.  Then:
\begin{align}
F_{\mu\nu}F^{\mu\nu}
&= \tilde{g}^{\mu\alpha}\tilde{g}^{\nu\beta}F_{\mu\nu}F_{\alpha\beta}
= 2\tilde{g}^{00}\tilde{g}^{ij}F_{0i}F_{0j}
  + \tilde{g}^{ij}\tilde{g}^{kl}F_{ik}F_{jl} \notag\\
&= 2\!\left(-\frac{e^{2\psi}}{c^2}\right)\delta^{ij}E_i E_j
  + \delta^{ij}\delta^{kl}\epsilon_{ikm}\epsilon_{jln}B_m B_n \notag\\
&= -\frac{2e^{2\psi}}{c^2}\,E^2 + 2B^2.
\end{align}
So:
\begin{equation}\label{eq:LEM-gordon-final}
\mathcal{L}_{\rm EM}^{\rm (Gordon)}
= -\frac{\sqrt{-\tilde{g}}}{4}\,F_{\mu\nu}F^{\mu\nu}
= -\frac{c\,e^{-\psi}}{4}\!\left(-\frac{2e^{2\psi}}{c^2}\,E^2 + 2B^2\right)
= \frac{e^{+\psi}}{2c}\,E^2 - \frac{c\,e^{-\psi}}{2}\,B^2.
\end{equation}

\paragraph{Variation with respect to $\psi$.}
\begin{equation}\label{eq:dLdpsi-gordon}
\frac{\partial\mathcal{L}_{\rm EM}^{\rm (Gordon)}}{\partial\psi}
= +\frac{e^{+\psi}}{2c}\,E^2 + \frac{c\,e^{-\psi}}{2}\,B^2.
\end{equation}

\begin{finding}[The Gordon Metric Does Produce an E--B Asymmetry!]\label{find:gordon-asym}
Despite the impedance-matched constitutive relations
$\varepsilon = \varepsilon_0 n$, $\mu = \mu_0 n$, the EM Lagrangian
on the Gordon metric (\ref{eq:LEM-gordon-final}) has different
$\psi$-dependences for $E$ and $B$:
\begin{equation}
\mathcal{L}_{\rm EM}^{\rm (Gordon)}
= \frac{e^{+\psi}}{2c}\,E^2 - \frac{c\,e^{-\psi}}{2}\,B^2.
\end{equation}
The $E^2$ term scales as $e^{+\psi}$ while the $B^2$ term scales as
$e^{-\psi}$.  This is an E--B asymmetry generated by the optical metric
itself.
\end{finding}

\begin{remark}[Resolution of the Apparent Contradiction]
Finding~\ref{find:gordon-asym} appears to contradict Theorem~\ref{thm:kappa-zero}.
The resolution is that the constitutive relations
$\varepsilon = \varepsilon_0 n$, $\mu = \mu_0 n$ refer to the
\textit{3D constitutive tensors} relating $\mathbf{D}$ to $\mathbf{E}$
and $\mathbf{H}$ to $\mathbf{B}$:
\begin{equation}
D^i = \varepsilon_{\rm eff}\,E^i = \varepsilon_0 n\,E^i, \qquad
H^i = \frac{1}{\mu_{\rm eff}}\,B^i = \frac{1}{\mu_0 n}\,B^i.
\end{equation}
The \textit{energy density} $u = \frac{1}{2}(\mathbf{D}\cdot\mathbf{E}
+ \mathbf{B}\cdot\mathbf{H})$ is
\begin{equation}
u = \frac{\varepsilon_0 n}{2}\,E^2 + \frac{1}{2\mu_0 n}\,B^2
= \frac{e^{+\psi}}{2}\!\left(\varepsilon_0 E^2 + \frac{B^2}{\mu_0 e^{2\psi}}\right).
\end{equation}
The 3D energy density has $\varepsilon_0 n E^2 + B^2/(\mu_0 n)$, which
after multiplied by the 3D volume element (which is unity on flat spatial
sections) gives contributions $e^{+\psi}E^2$ and $e^{-\psi}B^2$---exactly
matching Eq.~(\ref{eq:LEM-gordon-final}).

So the E--B asymmetry in the Lagrangian is \textbf{real and physical}:
it comes from the fact that $\varepsilon_{\rm eff} = \varepsilon_0 n$
affects $E$ while $1/\mu_{\rm eff} = 1/(\mu_0 n)$ affects $B$.  Even
though both $\varepsilon$ and $\mu$ increase by the same factor $n$,
the electric energy $\propto \varepsilon E^2 \propto nE^2$ increases
while the magnetic energy $\propto B^2/\mu \propto B^2/n$ decreases.
\end{remark}

\subsection{Extracting $\kappa$ from the Gordon Metric}

The constitutive relations are:
\begin{equation}
\varepsilon_{\rm eff} = \varepsilon_0\,e^{+\psi}, \qquad
\mu_{\rm eff} = \mu_0\,e^{+\psi}.
\end{equation}
The EM energy density (from the 3D perspective) is:
\begin{equation}
u_{\rm EM} = \frac{1}{2}\varepsilon_{\rm eff}\,E^2
+ \frac{B^2}{2\mu_{\rm eff}}
= \frac{\varepsilon_0}{2}\,e^{+\psi}\,E^2
+ \frac{1}{2\mu_0}\,e^{-\psi}\,B^2.
\end{equation}

Now compare with the $\kappa$ parametrization:
\begin{equation}
u_{\rm EM}^{(\kappa)} = \frac{\varepsilon_0}{2}\,n\,e^{+\kappa\psi}\,E^2
+ \frac{1}{2\mu_0}\,n\,e^{-\kappa\psi}\,B^2
= \frac{\varepsilon_0}{2}\,e^{(1+\kappa)\psi}\,E^2
+ \frac{1}{2\mu_0}\,e^{(1-\kappa)\psi}\,B^2.
\end{equation}

Matching:
\begin{align}
E\text{-sector:} &\quad e^{(1+\kappa)\psi} = e^{+\psi}
\quad\Rightarrow\quad \kappa = 0. \\[4pt]
B\text{-sector:} &\quad e^{(1-\kappa)\psi} = e^{-\psi}
\quad\Rightarrow\quad \kappa = 2.
\end{align}

These are \textbf{inconsistent} because the $\kappa$ parametrization
of Appendix~R uses a different convention.  In Appendix~R:
\begin{equation}
\varepsilon(\psi) = \varepsilon_0\,n\,e^{+\kappa\psi}, \qquad
\mu(\psi) = \mu_0\,n\,e^{-\kappa\psi}.
\end{equation}
The energy density is:
\begin{equation}
u = \frac{1}{2}\varepsilon E^2 + \frac{B^2}{2\mu}
= \frac{\varepsilon_0}{2}\,e^{(1+\kappa)\psi}\,E^2
+ \frac{1}{2\mu_0}\,e^{-(1-\kappa)\psi}\,B^2.
\end{equation}
Wait---$1/\mu = 1/(\mu_0 n\,e^{-\kappa\psi}) = e^{-(1-\kappa)\psi}/\mu_0$.

So matching the Gordon-metric result $e^{+\psi}E^2$ and $e^{-\psi}B^2$:
\begin{align}
E\text{-sector:} &\quad 1+\kappa = 1
\quad\Rightarrow\quad \kappa = 0. \\[4pt]
B\text{-sector:} &\quad -(1-\kappa) = -1
\quad\Rightarrow\quad \kappa = 0. \label{eq:kappa-zero-confirmed}
\end{align}

\begin{theorem}[$\kappa = 0$ from the Gordon Metric, Rigorously]\label{thm:kappa-zero-final}
The Gordon optical metric with $n = e^\psi$ and flat spatial sections
produces the EM Lagrangian
\begin{equation}
\mathcal{L}_{\rm EM} = \frac{e^{+\psi}}{2c}\,E^2
- \frac{c\,e^{-\psi}}{2}\,B^2.
\end{equation}
In the Appendix~R parametrization
$\varepsilon = \varepsilon_0\,n\,e^{+\kappa\psi}$,
$\mu = \mu_0\,n\,e^{-\kappa\psi}$, this corresponds to
$\kappa = 0$ for both sectors.

The E--B asymmetry in the Lagrangian ($e^{+\psi}$ vs $e^{-\psi}$) is
entirely accounted for by the standard constitutive relations
$\varepsilon = \varepsilon_0 n$, $\mu = \mu_0 n$; the $\kappa$
parameter measures the \textbf{residual} asymmetry beyond this,
which vanishes.
\end{theorem}

%=============================================================================
\section{The Auxiliary Metric and the Physical $\kappa$}
%=============================================================================

\subsection{The Auxiliary Metric as a Gauge-Sector Effective Metric}

The auxiliary metric $\hat{g}_{\mu\nu}$ used in the gauge emergence
derivations has $\hat{g}_{ij} = e^{+2\psi}\,\delta_{ij}$, doubling
the spatial exponent relative to the physical metric
$\tilde{g}_{ij} = e^{+\psi}\,\delta_{ij}$.

The EM Lagrangian on the auxiliary metric is:
\begin{equation}
\mathcal{L}_{\rm EM}^{(\rm aux)} = \frac{e^{+2\psi}}{2g^2 c}\,E^2
- \frac{c\,e^{-2\psi}}{2g^2}\,B^2.
\end{equation}

In the $\kappa$ parametrization with $n = e^\psi$:
\begin{align}
\varepsilon_{\rm aux} &= \varepsilon_0\,n\,e^{+\kappa_{\rm aux}\psi}
\propto e^{(1+\kappa_{\rm aux})\psi}
\stackrel{!}{=} e^{+2\psi}
\quad\Rightarrow\quad \kappa_{\rm aux} = 1, \\[4pt]
\frac{1}{\mu_{\rm aux}}
&= \frac{1}{\mu_0 n}\,e^{+\kappa_{\rm aux}\psi}
\propto e^{-(1-\kappa_{\rm aux})\psi}
\stackrel{!}{=} e^{+2\psi}
\quad\Rightarrow\quad \kappa_{\rm aux} = 3.
\end{align}
Still inconsistent.  The fundamental issue is that the auxiliary metric
does not preserve $v_{\rm ph} = c/n$ with $n = e^\psi$; it has
$v_{\rm ph,aux} = c\,e^{-2\psi}$, corresponding to $n_{\rm aux} = e^{2\psi}$.

\subsection{Correct Interpretation}

The $\kappa$ parametrization assumes the product
$\varepsilon\mu = \varepsilon_0\mu_0\,n^2$ is preserved.  For the
Gordon metric this holds ($\varepsilon_0 n \cdot \mu_0 n = \varepsilon_0\mu_0 n^2$).
For the auxiliary metric:
\begin{equation}
\varepsilon_{\rm aux}\mu_{\rm aux}
\propto e^{+2\psi}\cdot e^{-2\psi} = 1 \neq \varepsilon_0\mu_0 n^2.
\end{equation}
The product is $\varepsilon_0\mu_0$ (constant), not $\varepsilon_0\mu_0 n^2$.
The auxiliary metric gives $v_{\rm ph} = c$ (constant)---it does not
reproduce the DFD refractive index!

\begin{finding}[The Auxiliary Metric Cannot Be Parametrized by $\kappa$]
The $\kappa$ parametrization of Appendix~R is built on the assumption
that $v_{\rm ph} = c/n = c\,e^{-\psi}$.  The auxiliary metric
$\hat{g}_{\mu\nu}$ gives instead $v_{\rm ph} = c\,e^{-2\psi}$ (or
$v_{\rm ph} = c$, depending on normalization).  The $\kappa$ parameter
is therefore not directly applicable to the auxiliary-metric Lagrangian.
\end{finding}

%=============================================================================
\section{The Correct Route to Physical $\kappa$}
%=============================================================================

Given that $\kappa = 0$ from the minimal optical metric, and the
auxiliary-metric Lagrangian is a computational device, how does
$\kappa \neq 0$ arise physically?

\subsection{$\kappa$ as a Loop-Generated Parameter}

The tree-level result is $\kappa = 0$.  At one loop, gauge field
fluctuations on the optical-metric background generate effective
corrections to $\varepsilon$ and $\mu$ that are \textit{different}
for the electric and magnetic sectors.

The physical mechanism is:
\begin{enumerate}
\item Virtual charged particle loops (vacuum polarization) correct
$\varepsilon$, making it ``more permittive'' near charges.
\item The magnetic permeability receives no such correction at leading
order (vacuum magnetization is higher-order).
\item On a $\psi$-background, these corrections become $\psi$-dependent,
generating a $\kappa_{\rm eff} \neq 0$.
\end{enumerate}

The leading contribution is the Euler--Heisenberg correction, but
since you have asked for a derivation that does \textbf{not} invoke
Euler--Heisenberg, let us instead use the gauge emergence framework
directly.

\subsection{$\kappa$ from Gauge Emergence: The Correct Derivation}

The gauge emergence framework gives the gauge-$\psi$ Lagrangian
(Appendix~G) as:
\begin{equation}
\mathcal{L}_{\rm gauge-\psi} = \sum_r
\frac{c}{2g_r^2}\left[\frac{e^{2\psi}}{c^2}\,E_r^2 - e^{-2\psi}\,B_r^2\right].
\end{equation}
The physical EM Lagrangian on the Gordon metric is:
\begin{equation}
\mathcal{L}_{\rm EM}^{\rm (Gordon)} = \frac{1}{2c}
\left[e^{+\psi}\,E^2 - c^2\,e^{-\psi}\,B^2\right].
\end{equation}
(In appropriate units with $g_{\rm EM}^2 = 1$.)

The \textbf{ratio} of the auxiliary-metric coefficients to the
Gordon-metric coefficients defines the gauge-emergence correction:
\begin{align}
R_E &\equiv \frac{e^{+2\psi}}{e^{+\psi}} = e^{+\psi}, \\[4pt]
R_B &\equiv \frac{e^{-2\psi}}{e^{-\psi}} = e^{-\psi}.
\end{align}
The ratio $R_E/R_B = e^{+2\psi}$, which means the gauge-emergence
correction enhances $E$ relative to $B$ by $e^{+2\psi}$.

In the $\kappa$ language, this correction is
$e^{+\kappa_{\rm corr}\psi}$ for the $E$-sector and
$e^{-\kappa_{\rm corr}\psi}$ for the $B$-sector, with
$\kappa_{\rm corr} = 1$.

However, this correction must be \textbf{weighted by the gauge
coupling}.  The gauge emergence correction enters at order
$\alpha_{\rm eff}$:
\begin{equation}
\kappa_{\rm phys} = \kappa_{\rm tree} + \alpha_{\rm eff}\cdot\kappa_{\rm corr}
= 0 + \alpha_{\rm eff}\cdot 1 = \alpha_{\rm eff}.
\end{equation}

With $\alpha_{\rm eff} = \alpha/n_2^2 = \alpha/4$ (from Theorem~G.3):
\begin{equation}
\boxed{\kappa_{\rm phys} = \frac{\alpha}{4} \approx 1.82 \times 10^{-3}}
\end{equation}

\begin{finding}[$\kappa$ Is Not Independent of $\eta_c$]
The physical $\kappa$ arising from gauge emergence corrections to the
tree-level optical metric is precisely $\kappa = \alpha/4 = \eta_c$.
This identifies $\kappa$ with the EM-$\psi$ coupling threshold,
suggesting a deep connection: \textbf{the E--B split and the
nonlinear coupling threshold have a common origin in the $SU(2)$
frame stiffness.}
\end{finding}

%=============================================================================
\section{Summary and Open Questions}
%=============================================================================

\begin{tcolorbox}[colback=blue!5!white, colframe=blue!50!black, title=Summary of Results]
\begin{enumerate}
\item \textbf{Tree level:} The Gordon optical metric forces $\kappa = 0$.
The E and B sectors couple to $\psi$ symmetrically through the
impedance-matched constitutive relations
$\varepsilon = \varepsilon_0 n$, $\mu = \mu_0 n$.

\item \textbf{The ``E--B asymmetry'' in the Lagrangian} ($e^{+\psi}E^2$ vs
$e^{-\psi}B^2$) is \textit{not} a nonzero $\kappa$.  It is the standard
constitutive result: electric energy $\propto \varepsilon E^2 \propto nE^2$
increases with $\psi$, while magnetic energy $\propto B^2/\mu \propto B^2/n$
decreases.

\item \textbf{The auxiliary metric} used in gauge emergence gives a genuine
E--B asymmetry ($e^{+2\psi}E^2$ vs $e^{-2\psi}B^2$), but this metric
cannot be parametrized by $\kappa$ because it violates the
$v_{\rm ph} = c/n$ assumption.

\item \textbf{Physical $\kappa$} arises as a gauge-emergence correction
to the tree-level constitutive relations, weighted by $\alpha_{\rm eff}$:
\begin{equation}
\kappa = \frac{\alpha}{4} = \eta_c \approx 1.82 \times 10^{-3}.
\end{equation}

\item \textbf{This derivation uses no electroweak mixing angle and
no Euler--Heisenberg Lagrangian.}  It relies only on:
\begin{itemize}
\item The Gordon optical metric (DFD postulate)
\item The gauge emergence auxiliary metric (Appendix~G)
\item The frame stiffness $\kappa_2 = n_2\,\kappa_0$ (Theorem~F.6)
\item The effective coupling $\alpha_{\rm eff} = \alpha/n_2^2$ (Theorem~G.3)
\end{itemize}
\end{enumerate}
\end{tcolorbox}

\subsection{Open Questions}

\begin{enumerate}
\item \textbf{Is $\kappa = \alpha/4$ the full answer?}  Higher-order
corrections from the gauge emergence framework could modify this.
The $k_a$ derivation (Theorem~G.1) operates in the magnetically
dominated regime; does the electrically dominated regime give a
different effective $\kappa$?

\item \textbf{Is Appendix~R's $\kappa$ the same parameter?}
Appendix~R treats $\kappa$ as a free phenomenological parameter
constrained by experiment ($|\kappa| \lesssim 1$).  The value
$\kappa = \alpha/4 \approx 2 \times 10^{-3}$ is well within this
bound.  But Appendix~R's formalism allows arbitrary $\kappa$; the
gauge emergence result $\kappa = \alpha/4$ would be a prediction
to be tested.

\item \textbf{Falsification.}  The prediction $\kappa = \alpha/4$
can be tested via TE/TM polarization swaps in high-precision
cavity experiments (as described in Appendix~R, Section~R.6.4).
The predicted E--B split at order $10^{-3}$ in the exponent is
experimentally accessible.

\item \textbf{Relation to $k_\alpha$.}  The clock coupling
coefficient $k_\alpha = \alpha^2/(2\pi) \approx 8.5 \times 10^{-6}$
is much smaller than $\kappa = \alpha/4$.  These are different
parameters: $k_\alpha$ measures how atomic clocks respond to
$\psi$-variation, while $\kappa$ measures the differential
E--B response.  The ratio $\kappa/k_\alpha = \pi/(2\alpha) \approx 215$
is order $1/\alpha$, reflecting the tree-level vs.\ one-loop
distinction.
\end{enumerate}

\end{document}
